\documentclass[10pt,aps,pre,twocolumn,nopacs,final,letterpaper,superscriptaddress,longbibliography]{revtex4-2}
\usepackage[utf8]{inputenc}
\usepackage{amsmath,amsfonts,amssymb}
\usepackage{graphicx}
\usepackage{bm}
\usepackage{hyperref}
\usepackage[ruled]{algorithm2e}
\usepackage{xcolor}
\usepackage{quotes}
\definecolor{goodblue}{RGB}{0, 91, 187}
\urlstyle{same}
\hypersetup{
  colorlinks=true,
  allcolors=goodblue,
  urlcolor=goodblue,
  citecolor=goodblue,
  pdfborder={0 0 0},
  breaklinks=true,
}
%%% Macros
\newcommand{\boldA}{\bm{A}}
\newcommand{\boldQ}{\bm{Q}}
\newcommand{\boldX}{\bm{X}}
\newcommand{\bbN}{\mathbb{N}}
\newcommand{\boldx}{\bm{x}}
\newcommand{\boldc}{\bm{c}}
\newcommand{\boldq}{\bm{q}}
\newcommand{\boldone}{\bm{1}}
\newcommand{\ER}{Erd\H os-R\'enyi }

\begin{document}
\title{The complexity of contagion: nonlinearity and correlated heterogeneity}

\begin{abstract}

\end{abstract}

\maketitle

\section{Introduction} %LHD happy to write this to explain the origin of the project and the big research questions.

\begin{itemize}
    \item Models of epidemiological contagion are often modeled as simple contagion processes, based on a mechanistic understanding that a single infected individual can infect an un-infected individual.
    \item Talk about system-level biases emerging from uncertainty (St-Onge and H\'ebert-Dufresne)
\end{itemize}

\section{Results}

\subsection{The contagion model}

We consider a susceptible-infected-susceptible (SIS) contagion model, where susceptible nodes get infected via infected neighbors, and heal spontaneously. more specifically, we use a heterogeneous neighborhood-based contagion model, which is an extension of the model presented in Ref.~\cite{landry_reconstructing_2024}. Here, we allow each agent to have different rules of contagion spread.

More formally, given an undirected, simple, unweighted network $G = (V, E)$ with $N = |V|$ nodes and $M=|E|$ edges, for a node $i$, we define its neighborhood as $\eta_i = \{j \mid \forall (i, j)\in E\}$. The number of neighbors is, by definition, equal to the degree of node $i$.

We describe the discrete-time network SIS model as follows.
For each time $t=0, 1,\dots, T$, an infected node $i$ can heal spontaneously with probability $\gamma_i$, and be infected with probability $c_i(\nu_i)$, where $c_i\mapsto[0, 1]$ is the \textit{contagion function} for node $i$, and $\nu_i$ is the number of infected neighbors.
As in Ref.~\cite{landry_reconstructing_2024}, we write this function as a contagion vector $\boldc = [c_0, \dots, c_{N-1}]^T$  in nonparametric form, where $c_\nu$ denotes the probability of infection by $\nu$ infected neighbors and $N$ is the number of nodes.
The heterogeneity of the contagion function can capture, for instance, different immune responses characteristic of different ages, or different communities.
We can capture this by rewriting the contagion function as a hierarchical model $c(\nu, \theta) = c(\nu \mid \theta)P(\theta)$, where $\theta$ represents additional function parameters which make be contingent upon demographic characteristics.

Other things:
\begin{itemize}
    \item Our initial model is on a 2-community SBM with fixed mean degree, but one could (1) vary the relative community sizes and (2) add more than 2 communities. For example, in the three community case, you could arrange them in a line or a ring (both of which would be very interesting!)
    \item Would it be valuable to also look at contagion kernels conditioned on degree? Is this a special case of the $\theta$? I think it's different, because it's a structural property, although if the adjacency matrix is fixed, it's effectively the same.
\end{itemize}

\subsection{Contagion kernel inference}

Similar, to Ref.~\cite{landry_reconstructing_2024}, we use a Bayesian approach for inferring the contagion kernels. However, unlike Ref.~\cite{landry_reconstructing_2024}, we presume that the network itself is known and we must only infer the individual infection kernels.
There is a twist however: we infer the kernel using three different models: (1) a misspecified single-kernel model, (2) an individual kernel model, and (3) a hierarchical kernel model.

To-do:
\begin{itemize}
    \item The single-contagion-function model is pretty straightforward from Ref.~\cite{landry_reconstructing_2024}, and the individual kernel model is also straightforward (but needs much more data to faithfully estimate), but I haven't had a chance to write down the hierarchical model.
    \item Does it make sense to have some analysis on reconstructing the adjacency matrix from heterogeneous kernels? Or is that beyond the scope of this paper?
    \item It would be cool to get some analytical results relating the complexity of the contagion, the community structure and different contagion kernels.
\end{itemize}


\subsection{Potential model}

\begin{itemize}
    \item \textbf{Goal}: Understand variations in kernels. Can we get complex contagion from a mixture of kernels that look simple? 
    \item \textbf{Idea}: Each node gets its own kernel, but we pool across nodes. 
\end{itemize}

\textbf{Proposed model.}
Let $c_i(\nu)$ be the kernel for node $i$.
We know from previous work~\cite{landry_reconstructing_2024} that the likelihood of series of states $\bm{X}$ is proportional to
\begin{multline}
    \log P(\bm{X}\mid \bm{A}, \theta) \propto \sum_{i=1}^N \sum_{\nu=0}^{k_i} m_{i\nu} \log c_i(\nu)\\ +  \sum_{i=1}^N \sum_{\nu=0}^{k_i} n_{i\nu} \log \bigl[1 - c_i(\nu)\bigr]
\end{multline}
where $k_i$ is the degree of node $i$, and where $m_{i\nu}$ and $n_{i\nu}$ are the number of infection events occurring and failing to occur around node $i$ when it has $\nu$ infected neighbors.
The resulting matrices, $\bm{M}$ and $\bm{N}$, fully summarize the dynamics and are the only inputs required for inference.

In previous work, we made the simplifying assumption $c_i(\nu) = c(\nu)$ for all nodes $i=1,...,N$.
The proposed innovation for this paper is to model the heterogeneity of $\{c_i(\nu)\}_{i=1,...,N}$ more carefully.

We'll begin by putting this variable on a logit scale to allow for more flexible modeling
\begin{equation}
    \mathrm{logit}\ c_i(\nu) = \alpha_i(\nu),
\end{equation}
where $\alpha_i(\nu)\in\mathbb{R}$ is an intercept and now the main target of inference.
Then, we remove the assumption that every node has the same kernel and instead estimate the distribution of individual kernels using a hierarchical model where intercepts are random.

For example, we could model
\begin{equation}
    \alpha_i(\nu)\sim\mathrm{Normal}\bigl(\mu_\alpha(\nu), \sigma_\alpha\bigr)
\end{equation}
with further priors on the hyperparameters $\mu_\alpha$ and $\sigma_\alpha$.
(The normal distribution encourages clustering of the intercepts---a ridge regression of sorts. But the intercepts need not follow a normal distribution even if that's the prior. We can also use a different prior if this fits the data poorly.)
To complete the model, we need priors on the hyperparameters.
A simple prior on $\sigma_\alpha$ can be used, e.g., an exponential or Cauchy distribution with fixed hyperparameters.
The prior on $\mu_\alpha(\nu)$ can also be pretty generic, or we could enforce non-decreasing (average) intercepts with $\Delta\sim\mathrm{Dirichlet}(k_{\max})$ and
\begin{equation}
    \mu_\alpha(\nu) = \mathrm{logit}\left[\sum_{j=0}^{\nu} \Delta_j\right].
\end{equation}
(That is: We sample spacings between the average kernels, and the average is just a cumulative sum of the spacings. The logit transformation is necessary to work on the scale of the intercepts. )


The upshot is that this framework is easily extendable.
For example, we can bring in node features like degrees with
\begin{equation}
    \mathrm{logit}\ c_i(\nu) = \alpha_i(\nu) + \bm{x}_i^T\bm{\beta}
\end{equation}
where $\bm{x}_i = [x^{(i)}_1, ... x^{(i)}_q]$ is a vector of $q$ features for node $i$, and $\bm{\beta}$ are regression coefficients, to be inferred.
The rest of the model is unchanged.
(Note: We're essentially doing a logistic regression with random intercepts on the matrices $\bm{M}$ and $\bm{N}$.)


We can also write mixture versions, where nodes can be assigned one of many kernels.
This is the scenario where a node behaves differently depending on the situation, and we don't get to observe the kernel that was chosen.
We'd have many vectors of intercepts $\bm{\alpha}_1,...,\bm{\alpha}_K$, and
\begin{equation}
    \alpha_i(\nu) = \sum_{\ell=1}^K \pi_{i,\ell} \alpha_\ell (\nu)
\end{equation}
where $\bm{\pi}_i$ is drawn from a Dirichlet distribution and is different for every node.
(Again: the Dirichlet would have shared hyperparameters across nodes to enable pooling.)
In this case, a separate mixture for every node is probably too overparametrized; hence, the idea would be to introduce $K$ kernels. 
A specific node $i$ picks its kernel according to its personalized probability vector $\bm{\pi}_i$.
And this probability can depend on the feature of node $i$ if need be.
The model is otherwise the same.



\subsection{Case studies}

\subsubsection{Exploring empirical populations with SynthPops}

Big picture plan:
\begin{itemize}
    \item Collect age-mixing matrices from Mistry et al.~\cite{mistry_inferring_2021} and generate realizations of these matrices with \href{https://github.com/synthpops/synthpops}{SynthPops}.
    \item Get estimations of the infection probability for different age-classes.
    \item Measure contagion complexity for different strains and different cities?
    \item Is there an interesting question that can be answered by measuring the "complexity" of contagion across different empirical settings?
\end{itemize}

\section{Discussion}

Some ideas:
\begin{itemize}
    \item 
\end{itemize}


\section*{Data availability}

All code used to generate the results in this study are available on \href{https://github.com/nwlandry/heterogeneous-complex-network-reconstruction}{Github} \cite{zenodo}

\acknowledgments{N.W.L. acknowledges support from the University of Virginia Prominence-to-Preeminence (P2PE) STEM Targeted Initiatives Fund, SIF176A Contagion Science.}

\bibliography{references}
\end{document}
